\vspace{4pt plus 4pt minus 1pt}\noindent{\kbvsecfont\fontsize{11pt}{14.5pt}\selectfont\color{kbvaccentdark} (೬)\hspace{0.45em}ಮಹಾಸ್ರಗ್ಧರಾ ವೃತ್ತ :}\hspace{0.5em}ಇದು ನಾಲ್ಕು ಪಾದಗಳ ವೃತ್ತ. ಪ್ರತಿಯೊಂದು ಪಾದದಲ್ಲಿ ಏಳು ಗಣಗಳು ಮತ್ತು ಒಂದು ಗುರು ಸೇರಿ ಒಟ್ಟು ಇಪ್ಪತ್ತೆರಡು ಅಕ್ಷರಗಳಿರುವುವು.\par

\begin{kbvnote}
\noindent{\kbvsecfont ವಿನ್ಯಾಸ :}\hspace{0.5em}(ಸ + ತ + ತ + ನ + ಸ + ರ + ರ + ೧ಗು)
\end{kbvnote}

\vspace{6pt plus 5pt minus 1pt}\begin{center}{\kbvsecfont\fontsize{12.5pt}{16pt}\selectfont\color{kbvaccentdark} ಅಕ್ಷರ ಗಣಗಳ ವಿನ್ಯಾಸ ತೋರಿಸುವ ಪದಕ}\end{center}\vspace{3pt plus 3pt minus 1pt}

\vspace{5pt plus 4pt minus 1pt}
\noindent\begin{adjustbox}{max width=\textwidth}
{\renewcommand{\arraystretch}{1.30}\setlength{\tabcolsep}{4.5pt}
\begin{tabular}{l c c c c c c c c c c}
\toprule
\textbf{ವೃತ್ತ} & \textbf{ಗ1} & \textbf{ಗ2} & \textbf{ಗ3} & \textbf{ಗ4} & \textbf{ಗ5} & \textbf{ಗ6} & \textbf{ಗ7} & \textbf{ಲ} & \textbf{ಗು} & \textbf{ಒಟ್ಟು ಅಕ್ಷರಗಳು} \\
\midrule
ಉ.ಮಾ. & ಭ & ರ & ನ & ಭ & ಭ & ರ & .... & ೧ & ೧ & 20 \\
\hline
ಚಂ.ಮಾ. & ನ & ಜ & ಭ & ಜ & ಜ & ಜ & ರ & .... & .... & 21 \\
\hline
ಶಾ.ವಿ. & ಮ & ಸ & ಜ & ಸ & ತ & ತ & .... & .... & ೧ & 19 \\
\hline
ಮ.ವಿ. & ಸ & ಭ & ರ & ನ & ಮ & ಯ & .... & ೧ & ೧ & 20 \\
\hline
ಸ್ರಗ್ಧರಾ & ಮ & ರ & ಭ & ನ & ಯ & ಯ & ಯ & .... & .... & 21 \\
\hline
ಮ.ಸ್ರ. & ಸ & ತ & ತ & ನ & ಸ & ರ & ರ & .... & ೧ & 22 \\
\bottomrule
\end{tabular}}
\end{adjustbox}
\vspace{6pt plus 4pt minus 1pt}

\kbvexercise{ಅಭ್ಯಾಸ ಪತ್ರಿಕೆ : ೨೧}

\vspace{2.5pt plus 3pt minus 0.5pt}\noindent\hspace*{1.2em}\kbvnum{1.}\hspace{0.55em}ಛಂದಶ್ಶಾಸ್ತ್ರ ಎಂದರೇನು ?\par

\vspace{2.5pt plus 3pt minus 0.5pt}\noindent\hspace*{1.2em}\kbvnum{2.}\hspace{0.55em}“ಛಂದೋಂಬುಧಿ” ಎಂಬ ಗ್ರಂಥದ ಕರ್ತೃ ಯಾರು ?\par

\vspace{2.5pt plus 3pt minus 0.5pt}\noindent\hspace*{1.2em}\kbvnum{3.}\hspace{0.55em}“ತ್ರಿಪದಿ” ಎಂದರೇನು ?\par

\kbvsechead{4. ಇವುಗಳನ್ನು ಕುರಿತು ಟಿಪ್ಪಣಿ ಬರೆಯಿರಿ :}

\vspace{2.5pt plus 3pt minus 0.5pt}\noindent\hspace*{1.2em}\kbvnum{(೧)}\hspace{0.55em}ಯತಿ, (೨) ಗಣ, (೩) ಪ್ರಾಸ\par

\vspace{2.5pt plus 3pt minus 0.5pt}\noindent\hspace*{1.2em}\kbvnum{5.}\hspace{0.55em}ಇವುಗಳಲ್ಲಿ ಪ್ರತಿಯೊಂದಕ್ಕೆ ನಾಲ್ಕು ಉದಾಹರಣೆಗಳನ್ನು ಕೊಡಿ :\par

\vspace{2.5pt plus 3pt minus 0.5pt}\noindent\hspace*{1.2em}\kbvnum{(೧)}\hspace{0.55em}ಒಂದು ಮಾತ್ರಾಕಾಲದ ಅಕ್ಷರಗಳು\par

\vspace{2.5pt plus 3pt minus 0.5pt}\noindent\hspace*{1.2em}\kbvnum{(೨)}\hspace{0.55em}ಎರಡು ಮಾತ್ರಾಕಾಲದ ಅಕ್ಷರಗಳು\par
